% Part: first-order-logic
% Chapter: syntax-and-semantics
% Section: terms-formulas

\documentclass[../../../include/open-logic-section]{subfiles}

\begin{document}

\olfileid{fol}{syn}{frm}

\olsection{পদ ও \printtoken{P}{formula}}

একটি প্রথম-ক্রমের ভাষা~$\Lang L$ দেওয়া থাকলে, $\Lang L$-এর মৌলিক
শব্দভাণ্ডার থেকে গড়ে ওঠা রাশি আমরা সংজ্ঞায়িত করতে পারি। বিশেষত এর
মধ্যে পড়ে \emph{পদ} ও \emph{!!{formula}s}।

\begin{defn}[পদ]
\ollabel{defn:terms}
$\Lang L$-এর \emph{পদসমষ্টি}~$\Trm[L]$ নিচের নিয়মে আরোহীভাবে
সংজ্ঞায়িত:
\begin{enumerate}
\item প্রত্যেক !!{variable} একটি পদ।
\item $\Lang L$-এর প্রত্যেক !!{constant} একটি পদ।
\item $f$ যদি একটি $n$-স্থানীয় !!{function} হয় এবং $t_1$, \dots,
  $t_n$ পদ হয়, তাহলে $\Atom{f}{t_1, \ldots, t_n}$ একটি পদ।
\tagitem{limitClause}{
  আর কিছুই পদ নয়।}{}
\end{enumerate}
কোনো !!{variable}s নেই এমন পদকে \emph{বদ্ধ পদ} বলে।
\end{defn}

\begin{explain}
ভাষার বিবরণে এবং পদের সংজ্ঞায় !!{constant}s-কে আমরা একটি স্বতন্ত্র
সংকেতশ্রেণি হিসেবে রেখেছি, কিন্তু তার বদলে সেগুলিকে শূন্য-স্থানীয়
!!{function}s-এর মধ্যেও রাখা যেত। তাহলে পদের সংজ্ঞার দ্বিতীয় দফাটির
দরকার হতো না। শুধু বুঝতে হবে, $n = 0$ হলে
$\Atom{f}{t_1, \ldots, t_n}$ বলতে একা $f$-কেই বোঝায়।
\end{explain}

\begin{defn}[সূত্র]
\ollabel{defn:formulas}
$\Lang L$ ভাষার \emph{!!{formula}s-এর সেট}~$\Frm[L]$ নিচের নিয়মে
আরোহীভাবে সংজ্ঞায়িত:
\begin{enumerate}
\tagitem{prvFalse}{$\lfalse$ একটি পরমাণু !!{formula}।}{}

\tagitem{prvTrue}{$\ltrue$ একটি পরমাণু !!{formula}।}{}

\item $R$ যদি $\Lang L$-এর একটি $n$-স্থানীয় !!{predicate} হয় এবং
  $t_1$, \dots, $t_n$ যদি $\Lang L$-এর পদ হয়, তাহলে
  $\Atom{R}{t_1,\ldots, t_n}$ একটি পরমাণু !!{formula}।

\item $t_1$ ও $t_2$ যদি $\Lang L$-এর পদ হয়, তাহলে
  $\Atom{\eq}{t_1, t_2}$ একটি পরমাণু !!{formula}।

\tagitem{prvNot}{$!A$ যদি !!a{formula} হয়, তাহলে $\lnot !A$-ও
  !!a{formula}।}{}

\tagitem{prvAnd}{$!A$ ও $!B$ যদি !!{formula}s হয়, তাহলে $(!A \land
  !B)$-ও !!a{formula}।}{}

\tagitem{prvOr}{$!A$ ও $!B$ যদি !!{formula}s হয়, তাহলে $(!A \lor !B)$
  !!a{formula}।}{}

\tagitem{prvIf}{$!A$ ও $!B$ যদি !!{formula}s হয়, তাহলে $(!A \lif !B)$
  !!a{formula}।}{}

\tagitem{prvIff}{$!A$ ও $!B$ যদি !!{formula}s হয়, তাহলে $(!A \liff !B)$
  !!a{formula}।}{}

\tagitem{prvAll}{$!A$ যদি !!a{formula} এবং $x$ যদি !!a{variable} হয়,
  তাহলে $\lforall[x][!A]$ !!a{formula}।}{}

\tagitem{prvEx}{$!A$ যদি !!a{formula} এবং $x$ যদি !!a{variable} হয়,
  তাহলে $\lexists[x][!A]$ !!a{formula}।}{}

\tagitem{limitClause}{আর কিছুই !!a{formula} নয়।}{}
\end{enumerate}
\end{defn}

\begin{explain}
পদসমষ্টি ও !!{formula}s-এর সেটের সংজ্ঞা দুটি \emph{আরোহী সংজ্ঞা}।
মূলত আমরা অসীম সংখ্যক পর্যায়ে !!{formula}s-এর সেট নির্মাণ করি। প্রথম
পর্যায়ে সব পরমাণু সূত্রকে সূত্র বলে ঘোষণা করি; এটি সংজ্ঞার প্রথম কয়েকটি
ক্ষেত্রের, অর্থাৎ
\iftag{prvTrue}{$\ltrue$, }{}%
\iftag{prvFalse}{$\lfalse$, }{}%
$\Atom{R}{t_1,\dots,t_n}$ এবং $\Atom{\eq}{t_1,t_2}$-এর ক্ষেত্রের
অনুরূপ। তাই “পরমাণু !!{formula}” বলতে এই রূপের যেকোনো !!{formula}
বোঝায়।

সংজ্ঞার অন্য ক্ষেত্রগুলি ইতিমধ্যে নির্মিত !!{formula}s থেকে নতুন
!!{formula}s নির্মাণের নিয়ম দেয়। দ্বিতীয় পর্যায়ে এগুলি ব্যবহার করে
পরমাণু !!{formula}s থেকে !!{formula}s নির্মাণ করা যায়। তৃতীয় পর্যায়ে
পরমাণু সূত্রগুলি এবং দ্বিতীয় পর্যায়ে পাওয়া সূত্রগুলি থেকে নতুন সূত্র
নির্মাণ করি, এইভাবে ক্রমশ এগোই। কোনো একটি পর্যায়ে শেষ পর্যন্ত নির্মিত
হয় এমন সবকিছুই !!{formula}, আর কিছুই নয়।
\end{explain}

প্রথা অনুসারে $\eq$-কে তার যুক্তিগুলির মাঝে লিখি এবং বন্ধনী বাদ দিই:
$\eq[t_1][t_2]$ হলো $\Atom{\eq}{t_1,t_2}$-এর সংক্ষিপ্ত রূপ। তদুপরি,
$\lnot \Atom{\eq}{t_1,t_2}$-কে সংক্ষেপে $\eq/[t_1][t_2]$ লিখি।
$!B$, $!C$ থেকে কোনো দ্বি-স্থানীয় সংযোজক~$\ast$ ব্যবহার করে নির্মিত
$(!B \ast !C)$ সূত্রটি লেখার সময় আমরা প্রায়ই বাইরের একজোড়া বন্ধনী
বাদ দিয়ে শুধু~$!B \ast !C$ লিখব।

\begin{intro}
কিছু যুক্তিবিদ্যার বইয়ে
\iftag{prvEx}{$\lexists[x][!A]$ }{}%
\iftag{notprvEx,notprvAll}{}{ও }%
\iftag{prvAll}{$\lforall[x][!A]$ }{}%
যেন
\iftag{notprvEx,notprvAll}{!!a{formula}}{!!{formula}s}
হিসেবে গণ্য হয়, তার জন্য !!{variable}~$x$-কে $!A$-তে অবশ্যই সংঘটিত
হতে হয়। এই শর্তটি না রাখলে কোনো অসুবিধা হয় না, বরং কাজ সহজ হয়।
\end{intro}

\begin{tagblock}{defNot,defOr,defAnd,defIf,defIff,defTrue,defFalse,defEx,defAll}
\begin{defn}
সংজ্ঞায়িত অপারেটর দিয়ে নির্মিত সূত্রগুলি নিচের অর্থে বুঝতে হবে:

\begin{tagenumerate}{defTrue,defFalse,defNot,defOr,defAnd,defIf,defIff,defEx,defAll}

\tagitem{defTrue}{$\ltrue$ সংক্ষেপে বোঝায়
  \iftag{prvFalse}{$\lnot\lfalse$}{কোনো স্থির পরমাণু
    !!{formula}~$!A$-এর জন্য $(!A \lor \lnot !A)$}।}{}

\tagitem{defFalse}{$\lfalse$ সংক্ষেপে বোঝায়
  \iftag{prvTrue}{$\lnot\ltrue$}{কোনো স্থির পরমাণু
    !!{formula}~$!A$-এর জন্য $(!A \land \lnot !A)$}।}{}

\tagitem{defNot}{$\lnot !A$ সংক্ষেপে বোঝায় $!A \lif \lfalse$।}{}

\tagitem{defOr}{$!A \lor !B$ সংক্ষেপে বোঝায়
  \iftag{prvAnd}{$\lnot(\lnot !A \land \lnot !B)$}{$\lnot !A \lif
    !B$}।}{}

\tagitem{defAnd}{$!A \land !B$ সংক্ষেপে বোঝায়
  \iftag{prvOr}{$\lnot(\lnot !A \lor \lnot !B)$}{$\lnot (!A \lif
    \lnot !B)$}।}{}

\tagitem{defIf}{$!A \lif !B$ সংক্ষেপে বোঝায়
  \iftag{prvOr}{$\lnot !A \lor !B$}{$\lnot (!A \land \lnot !B)$}।}{}

\tagitem{defIff}{$!A \liff !B$ সংক্ষেপে বোঝায় $(!A \lif !B) \land (!B
  \lif !A)$।}{}

\tagitem{defAll}{$\lforall[x][!A]$ সংক্ষেপে বোঝায় $\lnot\lexists[x][\lnot !A]$।}{}

\tagitem{defEx}{$\lexists[x][!A]$ সংক্ষেপে বোঝায় $\lnot\lforall[x][\lnot !A]$।}{}
\end{tagenumerate}
% BN-SRC-088 TARGET-CORRECTION-BEGIN
\par\smallskip\noindent\textnormal{\small\textbf{উৎস-সংশোধন (BN-SRC-088):} হিমায়িত ইংরেজি উৎসে সংজ্ঞায়িত শর্তবচনের প্রথম বিকল্পটি $\lnot !A \lor !B)$; সেখানে কোনো মেলানো খোলা বন্ধনী ছাড়াই শেষে একটি অতিরিক্ত ডান বন্ধনী রয়েছে। বাংলা পাঠে সুগঠিত সংক্ষিপ্ত রূপ $\lnot !A \lor !B$ রাখা হয়েছে।} % BN-SRC-088 TARGET-CORRECTION-END
\end{defn}
\end{tagblock}

নির্দিষ্ট কোনো প্রয়োগের ভাষায় কাজ করলে আমরা প্রায়ই দ্বি-স্থানীয়
!!{predicate}s ও !!{function}s-কে তাদের সংশ্লিষ্ট পদগুলির মাঝে লিখব।
যেমন, পাটিগণিতের ভাষায় $t_1 < t_2$ ও $(t_1 + t_2)$ এবং সেটতত্ত্বের
ভাষায় $t_1 \in t_2$। পাটিগণিতের ভাষায় উত্তরসূরি ফলনটি প্রথামতো
তার যুক্তির \emph{পরেও} লেখা হয়:~$t'$। তবে বিধিবদ্ধভাবে এগুলি যথাক্রমে
$\Atom{\Obj A^2_0}{t_1, t_2}$, $\Obj f^2_0(t_1, t_2)$,
$\Atom{\Obj A^2_0}{t_1, t_2}$ ও $f^1_0(t)$-এর প্রচলিত সংক্ষিপ্ত রূপ।

\begin{defn}[সংকেতবিন্যাসগত অভিন্নতা]
$\ident$ সংকেতটি প্রতীকক্রমের মধ্যে সংকেতবিন্যাসগত অভিন্নতা প্রকাশ করে;
অর্থাৎ $!A \ident !B$ তখনই, যখন $!A$ ও $!B$ সমান দৈর্ঘ্যের প্রতীকক্রম
এবং প্রতিটি স্থানে একই সংকেত ধারণ করে।
\end{defn}

$\ident$ সংকেতের দুপাশে সংযুক্তকরণে পাওয়া প্রতীকক্রম থাকতে পারে। যেমন,
$!A \ident (!B \lor !C)$-এর অর্থ: $!A$ প্রতীকক্রমটি, এই ক্রমে একটি
খোলা বন্ধনী, $!B$ প্রতীকক্রম, $\lor$ সংকেত, $!C$ প্রতীকক্রম এবং একটি
বন্ধ বন্ধনী সংযুক্ত করে পাওয়া প্রতীকক্রমটিরই সমান। এটি সত্য হলে আমরা
জানি, $!A$-এর প্রথম সংকেত একটি খোলা বন্ধনী, $!A$-তে $!B$ একটি
উপপ্রতীকক্রম হিসেবে আছে (দ্বিতীয় সংকেত থেকে শুরু করে), তার পরে
$\lor$ আছে, ইত্যাদি।

পদ ও !!{formula}s আরোহী সংজ্ঞার মাধ্যমে মৌলিক উপাদান থেকে গড়ে ওঠে বলে
নিচের আরোহ-নীতিগুলি ব্যবহার করে আমরা তাদের সম্বন্ধে বিভিন্ন কথা প্রমাণ
করতে পারি।

\begin{lem}[\emph{পদের উপর আরোহের নীতি}]
    \ollabel{lem:trmind}
    ধরি $\Lang L$ একটি প্রথম-ক্রমের ভাষা।
    কোনো ধর্ম~$P$ এমন যে
    %
    \begin{enumerate}
        \item প্রত্যেক !!{variable}~$v$-এর ক্ষেত্রে ধর্মটি সত্য,
        %
        \item $\Lang L$-এর প্রত্যেক !!{constant}~$a$-এর ক্ষেত্রে ধর্মটি
        সত্য, এবং
        %
        \item $t_1$,~\dots, $t_n$-এর ক্ষেত্রে ধর্মটি সত্য এবং $f$ যদি
        $\Lang L$-এর একটি $n$-স্থানীয় !!{function} হয়, তাহলে
        $f(t_1,\dotsc,t_n)$-এর ক্ষেত্রেও ধর্মটি সত্য
    \end{enumerate}
    ($t_1$,~\dots, $t_n$-কে $\Lang{L}$-এর পদ ধরে),
    তাহলে $\Trm[L]$-এর প্রত্যেক পদের ক্ষেত্রেই $P$ সত্য।
\end{lem}

\begin{prob}
    \olref[fol][syn][frm]{lem:trmind} প্রমাণ করুন।
\end{prob}

\begin{lem}[\emph{!!{formula}s-এর উপর আরোহের নীতি}]
    \ollabel{thm:frmind}
    ধরি $\Lang L$ একটি প্রথম-ক্রমের ভাষা।
    কোনো ধর্ম~$P$ সব পরমাণু !!{formula}s-এর ক্ষেত্রে সত্য এবং এমন যে
    %
    \begin{enumerate}
        \tagitem{prvNot}{$!A$-এর ক্ষেত্রে সত্য হলে $\lnot !A$-এর
            ক্ষেত্রেও সত্য;}{}
        \tagitem{prvAnd}{$!A$ ও~$!B$-এর ক্ষেত্রে সত্য হলে
            $(!A \land !B)$-এর ক্ষেত্রেও সত্য;}{}
        \tagitem{prvOr}{$!A$ ও~$!B$-এর ক্ষেত্রে সত্য হলে
            $(!A \lor !B)$-এর ক্ষেত্রেও সত্য;}{}
        \tagitem{prvIf}{$!A$ ও~$!B$-এর ক্ষেত্রে সত্য হলে
            $(!A \lif !B)$-এর ক্ষেত্রেও সত্য;}{}
        \tagitem{prvIff}{$!A$ ও~$!B$-এর ক্ষেত্রে সত্য হলে
            $(!A \liff !B)$-এর ক্ষেত্রেও সত্য;}{}
        \tagitem{prvEx}{$!A$-এর ক্ষেত্রে সত্য হলে
            $\lexists[x][!A]$-এর ক্ষেত্রেও সত্য;}{}
        \tagitem{prvAll}{$!A$-এর ক্ষেত্রে সত্য হলে
            $\lforall[x][!A]$-এর ক্ষেত্রেও সত্য;}{}
  \end{enumerate}
  ($!A$ ও $!B$-কে $\Lang{L}$-এর !!{formula}s ধরে),
  তাহলে $\Frm[L]$-এর সব সূত্রের ক্ষেত্রে $P$ সত্য।
\end{lem}

\end{document}
